\documentclass[10pt,a4paper]{amsart}
\usepackage[margin=19mm]{geometry}
\usepackage{amsmath,amssymb,mathtools}
\usepackage[scr=rsfs,cal=euler]{mathalfa}
\usepackage{xfrac}
\usepackage[unicode,hidelinks,bookmarks=false]{hyperref}
\newcommand{\ie}{i.e.\ }
\newcommand{\cf}{cf.\ }
\newcommand{\resp}{respectively\ }
\newcommand{\Subsection}{\subsection}
\providecommand{\nobreakdash}{\nobreak\hbox{-}}
\setlength{\emergencystretch}{2em}
\multlinegap0pt
\usepackage{bm}
\usepackage{bbm}
\usepackage{dsfont}
\usepackage{xspace}
\usepackage{cancel}
\usepackage{mathtools}
\usepackage{amsthm}
\makeatletter
\@ifundefined{thm}{\newtheorem{thm}{Theorem}}{}
\@ifundefined{lem}{\newtheorem{lem}[thm]{Lemma}}{}
\makeatother
\def\BofH{\mathbb B(\mathcal H)}
\def\G{{\sf G}}
\def\X{\mathcal{X}}
\def\d{{\rm d}}
\newcommand{\A}{\mathfrak{A}}
\newcommand{\B}{\mathfrak{B}} 
\providecommand{\norm}[1]{\lVert#1\rVert}
\title[A note on finite-dimensional quotients and the problem of au]{A note on finite-dimensional quotients and the problem of automatic continuity for twisted convolution algebras}
\author{Felipe I. Flores}
\date{Source: arXiv:2410.12162v3 (CC BY 4.0)}
\begin{document}
\maketitle

\noindent\emph{Source and licence.} A note on finite-dimensional quotients and the problem of automatic continuity for twisted convolution algebras. Version arXiv:2410.12162v3 (CC BY 4.0), distributed under the Creative Commons Attribution 4.0 International licence, as recorded in the arXiv licence metadata for this exact version and shown on the abstract page https://arxiv.org/abs/2410.12162v3. Full rights notice: licenses/arxiv-2410.12162-CC-BY-4.0.txt.\par\smallskip

\noindent\textit{Editorial scope.} Counted unit: the Thm whose source label is \texttt{semisimplicity} in the article (source0.tex lines 371--373, source label \texttt{semisimplicity}), delivered together with the article's own complete original proof of it (source0.tex lines 374--419, one \texttt{proof} environment of 46 lines). No mathematics is written, repaired or re-derived: statement and proof are the article's own text line for line. Required context retained uncounted: exten (lines 333--339); easy (lines 354--363); inv (lines 359--361); dix (lines 367--369). Every definition, display and label of those lines is retained unchanged. Omitted from this package and not redistributed: 1--332, 340--353, 362--366, 370--370, 420--550. Nothing inside a retained excerpt is omitted or altered: every retained body is delivered exactly as the article's lines are written, so no omission receipt of any kind accompanies this package. Citations: the retained excerpts carry no citation command at all, so no bibliography entry of the article is transcribed and no reference is redistributed. Reference adaptation: none was needed and none was made; every reference inside the delivered unit resolves inside it, so no unresolved reference or citation is produced. Printed slips kept literal: no printed word, symbol, hypothesis or number of the counted statement or of its proof was altered. Layout and class: the article's own class is \texttt{article} with packages including appendix, xcolor, geometry, graphicx, subcaption, pdflscape, amssymb, ulem, hyperref, enumitem, mathrsfs, epstopdf, rotating, longtable; the wrapper is the lane's pinned \texttt{amsart} editorial wrapper at 10pt on A4, which restates the theorem-like environments the retained text needs and adds the article's own macro definitions that the retained text needs (\texttt{\textbackslash A}, \texttt{\textbackslash B}, \texttt{\textbackslash BofH}, \texttt{\textbackslash G}, \texttt{\textbackslash X}, \texttt{\textbackslash d}, \texttt{\textbackslash norm}), with extra wrapper packages (bm, bbm, dsfont, xspace, cancel, mathtools). The delivered numbering is the unit's own. Assets: no retained span contains a figure environment, a graphics-inclusion command, a PDF-page inclusion, a table or a TikZ picture; no image file is copied into this package, the package declares no asset dependency and no image file is redistributed with it. Licence: version arXiv:2410.12162v3 of this article is distributed under the Creative Commons Attribution 4.0 International licence, as recorded in the arXiv licence metadata for this exact version and shown on the abstract page https://arxiv.org/abs/2410.12162v3. A full rights notice is delivered at licenses/arxiv-2410.12162-CC-BY-4.0.txt. Characters: every retained excerpt is pure ASCII, so the delivered engine, XeLaTeX with its default Latin Modern text font, reports no missing character.
\begin{lem}\label{exten}
Let $\B$ be a Banach algebra, $\X$ be a Banach space, and let $\pi:\B\to \mathbb B(\X)$ be a contractive representation. Assume further that the following are true: \begin{itemize}
\item[(i)] $\B$ has a contractive approximate identity.
\item[(ii)] The representation $\pi$ is non-degenerate, that is, $\overline{\rm span}\{\pi(b)\xi\mid b\in\B, \xi\in \X\}=\X$.
\end{itemize}
Then there exists a unique contractive unital representation $\widetilde\pi: \mathcal M(\B)\to \mathbb B(\B)$, such that $\widetilde\pi\circ\iota_\B=\pi$.
\end{lem}

\begin{lem}\label{easy} The following statements are true.
\begin{enumerate}
\item[(i)] $\Gamma_{\G,\A}$ is a group.
\item[(ii)] Every $m_{u,y}\in \Gamma_{\G,\A}$ is unitary and has norm $1$.
\item[(iii)] Every multiplier of the form $m_{a,y}$ can be written as a linear combination of $4$ elements in $\Gamma_{\G,\A}$.
\item[(iv)] The adjoint of $m_{a,y}$ satisfies the formula \begin{equation}\label{inv}
m_{a,y}^{*}=m_{\omega(y^{-1},y)^*\alpha_{y^{-1}}(a^*),y^{-1}},
\end{equation} for all $a\in\mathcal{M}(\A), y\in\G$.
\end{enumerate}
\end{lem}

\begin{lem}\label{dix}
Let $\X$ be a finite-dimensional Hilbert space and $V\subset {\rm GL}(\X)$ be a subgroup such that $\sup_{v\in V}\norm{v}_{\BofH}<\infty$. Then there exists a positive and invertible linear transformation $T\in {\rm GL}(\X)$ such that $TvT^{-1}\in \mathcal U(\X)$, for all $v\in V$.
\end{lem}

\begin{thm}\label{semisimplicity}
Let $(\G,\alpha,\omega,\A)$ be a twisted action. If $I\subset \B=L^1_{\alpha,\omega}(\G,\A)$ is a closed, finite-codimensional two-sided ideal, then $I$ is automatically self-adjoint and the quotient algebra $\B/I$ is semisimple.
\end{thm}

\begin{proof}
Since $I$ is closed and finite-codimensional, $\mathcal X=\B/I$ is a finite-dimensional Banach space. Let $\langle\cdot,\cdot\rangle$ be any inner product; since it is finite-dimensional, $\X$ is a Hilbert space with respect to this inner product.

We denote by $\pi:\B\to\mathbb B(\mathcal X)$ the induced representation on the quotient, that is, 
$$
\pi(\Phi)(\Psi+I)=\Phi*\Psi+I,
$$
for all $\Phi,\Psi\in \B$. This representation is contractive and non-degenerate, so, due to Lemma \ref{exten} and abusing the notation, $\pi$ extends to $\mathcal M(\B)$ and, hence, the operators $\pi(m_{a,y})\in\mathbb B(\X)$ are well-defined and uniformly bounded. In fact, it is not difficult to notice that they satisfy the identity
$$
\pi(m_{a,y})(\Psi+I)=m_{a,y}(\Psi)+I,\quad\quad\text{ for all }\Psi\in\B.
$$
For this reason, one observes that
\begin{equation}\label{desinte}
\pi(\Phi)(\Psi+I)=\int_\G \pi(m_{\Phi(y),y})(\Psi+I)\d y=\int_\G m_{\Phi(y),y}(\Psi)\d y +I.
\end{equation}
Now, we note that $V=\{\pi(m)\}_{m\in\Gamma_{\G,\A}}$ satisfies all the conditions of Lemma \ref{dix} and, therefore, there must exist a positive and invertible operator $T\in {\rm GL}(\X)$ such that $T\pi(m)T^{-1}\in \mathcal U(\X)$, for all ${m\in\Gamma_{\G,\A}}$.

We then define the representation $\pi':\B\to\mathbb B(\X)$ given by $\pi'(\Phi)=T\pi(\Phi)T^{-1}$ and we will now prove that it is a $^*$-representation and that ${\rm Ker}\,\pi'=I$, from which it will follow that $\B/I$ is $^*$-isomorphic to $\pi'(\B)$, which is a $C^*$-algebra, and therefore we will have shown that $\B/I$ is semisimple.

Indeed, note that ${\rm Ker}\,\pi'={\rm Ker}\,\pi$. Now, let $\Phi\in{\rm Ker}\,\pi$ and let $\Psi_j\in\B$ be some approximate bounded identity of $\B$. We note that
$$
I=\lim_j\pi(\Phi)(\Psi_j+I)=\lim_j \Phi*\Psi_j+I=\Phi+I,
$$ 
so $\Phi\in I$. This proves that ${\rm Ker}\,\pi'=I$.

Now let's see that $\pi'$ is a $^*$-representation. Indeed, if $m\in\Gamma_{\G,\A}$ and $\xi,\eta\in\X$, then one has 
\begin{align*}
    \langle \pi'(m)\xi,\eta\rangle&=\langle \xi,\pi'(m)^{*}\eta\rangle \\
    &=\langle \xi,\pi'(m)^{-1}\eta\rangle \\
    &=\langle \xi,\pi'(m^{-1})\eta\rangle=\langle \xi,\pi'(m^*)\eta\rangle.
\end{align*}

But, remembering that every $m_{a,y}$ can be written as a linear combination of $4$ elements in $\Gamma_{\G,\A}$ (point \emph{(iii)} of Lemma \ref{easy}), we see that
$$
\pi'(m_{a,y})^*=\pi'(m^*_{a,y}),\quad\text{ for every }a\in\mathcal{M}(\A),y\in \G.
$$
And, consequently, for $\Phi\in \B,\xi\in\X$, and using the equality \eqref{desinte}, one observes that
\begin{align*}
    \pi'(\Phi^*)\xi=T\pi(\Phi^*)T^{-1}\xi&=T\int_\G \pi\big( m_{\Phi^*(y),y}\big)T^{-1}\xi\,\d y \\
    &= T\int_\G \Delta(y^{-1})\pi\big(m_{\omega(y,y^{-1})^*\alpha_y(\Phi(y^{-1})^*),y}\big)T^{-1}\xi\,\d y \\
    &=T\int_\G \pi\big(m_{\omega(y^{-1},y)^*\alpha_{y^{-1}}(\Phi(y)^*),y^{-1}}\big)T^{-1}\xi\,\d y \\
    &\overset{\eqref{inv}}{=}\int_\G T\pi\big(m_{\Phi(y),y}^*\big)T^{-1}\xi\,\d y \\
    &=\int_\G \pi'(m_{\Phi(y),y})^*\xi\,\d y=\pi'(\Phi)^*\xi,
\end{align*}
 which finishes the proof.
\end{proof}

\end{document}
